\subsection{States and Resolutions}

We select five states that together cover the major dimensions of
redistricting difficulty: North Carolina (NC, $k=14$), Texas (TX, $k=38$),
Wisconsin (WI, $k=8$), Florida (FL, $k=28$), and Virginia (VA, $k=11$).
North Carolina is chosen because it is the most frequently litigated
redistricting state of the modern era and has been central to prior bisect
portfolio papers.  Texas provides the largest single-state test ($k=38$) and
has significant minority populations.  Wisconsin offers a small-$k$ case
in a highly competitive partisan environment.  Florida and Virginia provide
Sunbelt diversity and mid-size district counts.

For each state we run bisect at three resolution levels:

\begin{description}
  \item[County (coarse).] County-equivalent boundaries from the 2020
    TIGER/Line files.  North Carolina has 100 counties ($n=100$, $k=14$),
    yielding a graph with $\sim$500 edges.  This is the sparsest possible
    resolution for which contiguous districts are achievable.
  \item[Tract (standard).] 2020 census tracts, the default for the bisect
    portfolio.  NC: $n \approx 2{,}200$; TX: $n \approx 5{,}300$; WI:
    $n \approx 1{,}900$; FL: $n \approx 4{,}200$; VA: $n \approx 2{,}000$.
  \item[Block group (fine).] 2020 block groups.  NC: $n \approx 5{,}000$;
    TX: $n \approx 14{,}000$; WI: $n \approx 4{,}700$; FL: $n \approx
    11{,}000$; VA: $n \approx 5{,}200$.
\end{description}

\subsection{Algorithm Parameters}

To isolate the resolution effect, all parameters are held fixed across
resolutions:

\begin{itemize}
  \item Population balance tolerance: $\pm0.5\%$ of ideal district
        population.
  \item Edge weights: geographic (shared boundary length, unscaled).
  \item Search strategy: \texttt{single} (seed 42 for reproducibility).
  \item METIS engine: \texttt{c-ffi} (bundled C METIS 5.1.0).
  \item Census year: 2020.
\end{itemize}

The \texttt{single} search strategy runs one METIS call per bisection node
and selects the first feasible partition.  This is the simplest possible
search mode, appropriate for a sensitivity study where the goal is
understanding variation across resolutions rather than optimising within a
resolution.

\textbf{Single-run caveat.}\textsuperscript{\dag}  All reported empirical
results in this paper derive from single runs (one seed, one METIS call per
bisection node).  Compactness and seat-count figures carry the
\textsuperscript{\dag} mark throughout to indicate they are not ensemble
averages.  The purpose of the paper is to demonstrate \emph{robustness across
resolutions}, not to characterise the full distribution of outcomes within
any resolution.  Single-run results are appropriate for cross-resolution
comparison because METIS is deterministic given a fixed seed.

\subsection{Data Sources}

\begin{description}
  \item[Population.] 2020 Census PL 94-171 Summary File, total population
    (P1 table).  Available at census tract, block group, and county level
    from Census Bureau API.
  \item[Geography.] 2020 TIGER/Line shapefiles, all three resolution levels,
    downloaded via \texttt{bisect fetch --year 2020}.
  \item[Election data.] 2020 presidential election results disaggregated to
    census tract level from the MIT Election Data Science Lab
    \citep{MEDSL2020}, then re-aggregated to block-group and county level
    by population-proportional allocation.
\end{description}

\subsection{Metrics}

For each state-resolution combination we compute:

\paragraph{Polsby-Popper compactness (PP).}
\begin{equation}
  PP = \frac{4\pi A}{P^2}
\end{equation}
where $A$ is district area and $P$ is district perimeter, computed on
dissolved district polygons in WGS84 coordinates.  We report the mean PP
across all $k$ districts in each state.

\paragraph{D-seat count.}
The number of districts in which the Democratic presidential vote share
exceeds 50\%, using 2020 presidential returns as the electoral benchmark.
This metric is deliberately simple---it uses a single election rather than
an average---consistent with the single-run framing.

\paragraph{Resolution stability index (RSI).}
To summarise cross-resolution variation, we define:
\begin{equation}
  \text{RSI}(s) = \frac{\max_r PP(s, r) - \min_r PP(s, r)}{\text{mean}_r PP(s, r)}
\end{equation}
where $r$ indexes resolution (county, tract, block group) and $s$ indexes
state.  An RSI $< 0.02$ means the range of PP across resolutions is less
than 2\% of the mean---our primary stability criterion.

\subsection{Comparison Baseline}

We also compare block-group-level results to the 2020 enacted congressional
district maps (where available from the NCSL districting data portal) to
situate the algorithmic results in the context of real-world practice.
This comparison is secondary to the cross-resolution analysis.

\subsection{Computational Environment}

All runs use the \texttt{bisect} binary (Rust, version matching the
\texttt{main} branch at submission time) on a Windows 11 workstation with 32
GB RAM.  Block-group-level runs for large states (TX, FL) require
approximately 6 GB peak memory and complete within 4 minutes on 8 threads.
County-level runs for all states complete in under 10 seconds.
